%%% rest.tex — die Elemente, die weder raster.tex noch kaesten.tex zeigen:
%%% Kapitelanfang mit Bild, Kreis- und Freiformmaske, Inhaltsverzeichnis,
%%% Rueckseite. Ohne Seitenhintergrund, damit ein Lauf kurz ist.
%%%
%%%   xelatex -output-directory=beispiel beispiel/rest.tex   (dreimal)
%%%   python3 werkzeuge/nachmessen.py beispiel/rest.pdf --text

\documentclass[raster,ohnehintergrund]{dsa5latex}

\begin{document}

\dsaInhalt

\dsakapitel[grafiken/aventurienkarte]{Kapitel mit Bild}

\dsaKapitelseiteEnde

Text nach der Kapitelanfangsseite. Er muss auf dem Grundlinienraster sitzen
und darf nicht hinter dem Kapitelbild liegen.

\dsaabschnitt{Masken}

\dsaBildRaster{7}{\dsaBildKreis{25mm}{grafiken/aventurienkarte}}

Text nach dem Kreisbild.

\dsaBildRaster{6}{\dsaBildForm{30mm}{20mm}{(0,0) (1,0.6) (1.6,-0.3) (0.6,-1)}{grafiken/aventurienkarte}}

Text nach der Freiform.

\dsaabschnitt{Bild in Spaltenbreite}

\dsaBildSpalte{grafiken/aventurienkarte}{10}

Darunter laeuft der Text weiter und muss wieder auf der Grundlinie sitzen,
weil das Bild eine ganze Zahl Rastereinheiten belegt. Noch ein Satz, damit
mehrere Zeilen entstehen und die Lage zu messen ist.

\dsaUmschlagHinten{grafiken/aventurienkarte}{Ein Beispiel}{%
  Hier steht der Klappentext, der auf der Rueckseite im Rahmen sitzen soll.

  \dsaZusammenfassung{ein Beispiel in wenigen Worten}{Ermittlung}{keine}%
    {irgendwo}{jederzeit}

  \dsaAnforderungen{1}{4}{2}{1}
}

\end{document}
